% ============================================================
% 付録G 記号・用語・定義の対応表
% ------------------------------------------------------------
%
% ============================================================

\chapter*{付録 G：記号・用語・定義の対応表}
\addcontentsline{toc}{chapter}{付録 G：記号・用語・定義の対応表}

本論文で用いる記号（$\mu, D, \sigma, S$）は，
理論モデル（Merton 型）と実証モデル（Ordered Logit）で
同じ構造的役割を担うが，
観測可能性の制約により，実証では理論上の定義を
財務データに写像した proxy を用いて推定している。
そのため，記号と用語の対応は文脈に応じて変化し，
単一の用語に固定されるものではない。
以下に，本論文での記号・用語・proxy の対応関係を整理する。

\begin{table}[h]
\centering
\renewcommand{\arraystretch}{1.35}
\begin{tabular}{p{1.5cm} p{3.8cm} p{3.8cm} p{5.0cm}}
\hline
\textbf{記号} &
\textbf{理論モデルでの定義} &
\textbf{実証モデルでの用語} &
\textbf{データ構築での proxy（実際に用いる変数）} \\
\hline

$\mu$ &
企業価値の期待成長率（収益性）\newline
log-diff による drift &
収益性（企業価値の期待成長率） &
総資産成長率（asset\_growth）\newline
理論では企業価値の成長率を用いるが，実証では観測可能な財務データに写像するため \\

$D$ &
負債負担（企業価値に対する負債の重さ） &
負債比率（財務レバレッジ） &
負債比率（liabilities / assets）\newline
理論では負債の重さとして潜在変数 $y^*$ を押し下げる構造的要因であり，
実証では貸借対照表の負債比率で代理するため \\

$\sigma$ &
企業価値のボラティリティ（ショック感応度） &
ショック感応度 &
営業利益変動率，CF volatility などの proxy\newline
企業価値のボラティリティは観測できないため，財務データの変動率で代理するため \\

$S$ &
支援能力（親会社・制度による緩衝材） &
支援能力（倒産閾値を押し下げる要因） &
日本：$S=0$（観測不可）\newline
英国：親会社の財務余力・資本注入履歴・制度的緩衝材（CfD, Capacity Market, SoLR）\newline
支援能力は国によって観測可能性が異なるため \\

$y^*$ &
潜在変数（倒産ステージを規定する連続量） &
潜在変数（Ordered Logit の線形予測子） &
$\mu, D, \sigma, S$ の線形結合\newline
理論と実証で同じ構造だが，実証では推定値として扱われるため \\

$\tau_1, \tau_2$ &
cutpoint（閾値） &
cutpoint（stage 0→1，1→2 の境界） &
推定される cutpoint\newline
理論では境界値，実証では推定されるパラメータ \\

\hline
\end{tabular}
\end{table}

